%%
%% Berliner Hochschule für Technik --  Abschlussarbeit
%%
%% Kapitel 2
%%
%%

\chapter{Rechtliche Rahmenbedingungen}
\section{Teilnahme}
Für die Teilnahme an der CueLens-Studie werden Teilnahmevereinbarungen mit Privatpersonen getroffen. Eine Teilnahme ist nur möglich, wenn alle Bedingungen erfüllt sind: Teilnehmende müssen deutsch oder englisch lesen können, über ein Android-Endgerät mit Internetzugang verfügen, dauerhaft in Deutschland oder der Schweiz leben, zwischen 30 und 65 Jahre alt sein und mindestens zehn Zigaretten pro Tag rauchen. Außerdem darf jede Person nur einmal teilnehmen, und es müssen 20 Selbstberichte abgegeben werden. 

\subsection{Prüfung der Teilnahmebedingungen}
Während das Sprachverständnis bereits die Anmeldung und das Endgerät die Teilnahme an der Studie überhaupt erst ermöglichen, ist der gewöhnliche Aufenthaltsort für die Teilnehmenden vor allem in Fragen wichtig, die den Schutz ihrer Daten betreffen. Das Alter und der Rauchstatus werden mittels Selbstauskunft überprüft. Eine vorsätzlich falsche Selbstauskunft ist gemäß § 123 BGB beziehungsweise Art. 28 OR vertragswidrig. Insofern auffällige Selbstberichte müssen nicht von der Auswertung ausgeschlossen, können aber als Ausreißer interpretiert werden. Einmalige Teilnahme kann praktikabel anhand von Namen und Bankverbindungen festgestellt werden. Die Anzahl der Selbstberichte pro App-Installation kann datenschutzgerecht gespeichert und zur Verfizierung herangezogen werden. 

Die Ausweisprüfung und ein klinischer Test des Kohlenmonoxid-Gehalts im Atem oder des Cotiningehalts im Blut wären effektiver, um Teilnahme, Alter beziehungsweise Rauchstatus zu überprüfen. Im Rahmen dieser Masterarbeit sind diese Maßnahmen nicht verhältnismäßig und wären zum Teil auch unzulässig.
 
\subsection{Teilnahmeprämie}
Nach abgeschlossener Teilnahme an der Studie wird eine vereinbarte Aufwandsentschädigung an die Teilnehmenden gezahlt.

\section{Datenschutz}
Für die CueLens-Studie werden personenbezogene, gesundheitsbezogene Daten erhoben. Diese Daten werden auf Servern in der EU gespeichert und in die Schweiz übertragen. Die Gesetze der beiden Rechtsordnungen sind ähnlich, und die grenzüberschreitende Übertragung ist unproblematisch, wenn die Informationspflichten gegenüber den Teilnehmenden eingehalten werden.

Der Artikel 1 der Verordnung (EU) 2016/679 (Datenschutz-Grundverordnung, DSGVO) nennt den Schutz der Grundrechte natürlicher Personen bei der Verarbeitung personenbezogener Daten sowie den Verkehr personenbezogener Daten in der Europäischen Union als Gegenstand und Ziele. Das schweizerische Bundesgesetz über den Datenschutz (Datenschutzgesetz, DSG) nennt in Artikel 1 den Zweck, die Grundrechte von natürlichen Personen, über die Personendaten bearbeitet werden, zu schützen. Der wichtigste konzeptionelle Unterschied zwischen beiden Gesetzen ist das grundsätzliche Verbot der Verarbeitung personenbezogener Daten, außer bei Vorliegen bestimmter Bedingungen, in der DSGVO, und die ausdrückliche Erlaubnis bei Vorliegen bestimmter Bedingungen im DSG. Die Bedinungen sind größtenteils gleich, sodass beide Gesetze den gleichen Effekt haben.

Das vorliegende Dokument enthält keine personenbezogene Daten von Studienteilnehmern.

\subsection{Zweck der Erhebung}
In Art. 5 DSGVO und Art. 6 DSG ist festgelegt, dass nur die für einen eindeutigen Zweck erforderlichen minimalen Datenmengen erhoben werden dürfen, dass die betroffenen Personen die Erhebung nachvollziehen können müssen, und dass die personenbezogenen Daten nur für die Dauer der zweckentsprechenden Bearbeitung gespeichert werden dürfen. 

Für die CueLens-Studie werden Name, E-Mail-Adresse, Alter, ungefähre Anzahl gerauchter Zigaretten pro Tag, 20 Selbstberichte zum Rauchverlangen und Zahlungsverbindungen der Teilnehmenden benötigt. 

\subsection{Einwilligung}
Weiterhin müssen die Teilnehmenden vor der Teilnahme in die Verwendung ihrer personenbezogenen Daten (Art. 6 DSGVO, Art. 6 DSG) beziehungsweise ihrer Gesundheitsdaten (Art. 9 DSGVO, Art. 6 DSG) einwilligen. Sie haben nach Art. 7 DSGVO beziehungsweise Art. 30 DSG das Recht, ihre Einwilligung zu widerrufen. Die bis zum Zeitpunkt eines Widerrufs erhobenen Daten dürfen aber weiterhin verwendet werden. 

Der Nachweis der Einwilligung wird gemäß Art. 7 DSGVO beziehungsweise Art. 6 DSG aufbewahrt. Dabei wird bei deutschen Teilnehmenden die dreijährige regelmäßige Verjährungsfirst nach § 195 BGB und bei Schweizer Teilnehmenden die Frist von zehn Jahren nach Art. 127 OR berücksichtigt.

\subsection{Informationspflichten}
In Art. 13 DSGVO und Art. 19 DSG ist vorgesehen, dass betroffenen Personen zum Zeitpunkt der Erhebung ihrer personenbezogenen Daten der Name und die Kontaktdaten des Verantwortlichen mitgeteilt werden. In der CueLens-Studie heißt der Verantwortliche Stefan Berger und ist unter der E-Mail-Adresse datenschutz@each-and-every.de kontaktierbar. Weiterhin müssen den Teilnehmenden der Zweck und die Rechtsgrundlage der Verarbeitung, die Kategorien der personenbezogenen Daten, die Empfänger der Daten und die Übertragung der Daten in die Schweiz beziehungsweise nach Deutschland mit Angemessenheitsbeschluss gemäß Art. 45 DSGVO beziehungsweise Art. 16 DSG mitgeteilt werden. 

Im Anhang \ref{Datenschutzerklärung} befindet sich die Datenschutzerklärung, die für die Teilnehmenden erstellt wurde.

\subsection{Technische Maßnahmen}
Die technischen Komponenten, die für die Studie verwendet werden, müssen nach Art. 25 und 32 DSGVO beziehungsweise Art. 7 und 8 DSG so ausgelegt und eingestellt werden, dass die Datenschutzgrundsätze eingehalten werden. Der Austausch von Daten geschieht ausschließlich über verschlüsselte HTTPS- oder SFPT-Verbindungen. Bei Zugriffsberechtigungen für Quelltext sowie Daten und Funktionen der App und des Servers wird das Prinzip der geringsten Privilegien angewandt. Die App überträgt maximal 20 Selbstberichte, und die Selbstberichte werden anonymisiert. Gegen missbräuchliche Serveranfragen werden verfügbare Dienste zur App-Integrität genutzt.

Alle Teilnahmedaten werden über den Webspace-Anbieter united-domains.de auf einem Anwendungsserver in Deutschland gespeichert, und es wird eine Sicherungskopie beim Webspace-Anbieter infomaniak.com in der Schweiz gepflegt. Es wurde mit united-domains GmbH und Infomania\-k Network AG jeweils eine Vereinbarung über die Verarbeitung von Daten gemäß Art.~28 DSGVO beziehungsweise Art. 9 DSG getroffen. 

\subsection{Datenkategorien}
Studienergebnisse werden ausschließlich aus anonymisierten Daten gewonnen. Kontakt- und Zahlungsinformationen werden für die Überprüfung der Teilnahme und Abrechnung der Aufwandsentschädigung benötigt. 

Die in diesem Kapitel genannten Grundsätze und Vorgehensweisen werden für die Dauer dieser Arbeit, insbesondere für die Appentwicklung, die Planung und Durchführung der App-Studie und der Auswertung der Berichte, eingehalten.

\section{Urheber- und Markenrecht}
Bei der Präsentation von Auslöserreizen und der Klassifizierung von Bildern müssen die Rechte der Eigentümer der Dateien, der abgebildeten Personen und der Inhaber der abgebildeten Marken beachtet werden. Die einfachste Methode, mit der diese Rechte gewahrt bleiben, ist die Verwendung selbst erstellter Bilder ohne erkennbare Personen und Marken.

\subsection{Auslöserreize}
Die Bilder für Cue-Matching und Cue-Labeling werden entweder mit einer Kamera selbst erstellt, durch KI generiert oder modelliert. Dabei wird darauf geachtet, dass die Bilder frei von urheberrechtlich und markenrechtlich geschützten Inhalten und erkennbaren Personen sind.

\subsection{Bildklassifizierung}
Während der Nutzung der App können eigene Bilddateien und Kamerabilder verwendet werden. Bilder aus diesen Quellen werden ausschließlich innerhalb der App für die Klassifizierung und direkt anschließendes Cue-Matching oder Cue-Labeling verwendet. Danach werden sie verworfen und nicht gespeichert. Insbesondere werden durch die App keine Bilder an jegliche Empfänger übermittelt.

\subsection{"`CueLens"'}
  Der Name der Studie ist CueLens. CueLens ist keine Marke. In deutschen (DPMA), schweize\-rischen (IGE), europäischen (EUIPO) und globalen (WIPO) Markenregistern sowie im Google Play Store wurden zum Suchbegriff "`CueLens"' keine bestehende Einträge gefunden.

